%% =====================================================================
%%  第 6 章　“没事”
%%  源文件：06-没事.md
%% =====================================================================
\chapter{“没事”}

\applicable{你问对方怎么了，对方回了两个字。}

\begin{guideblock}
\textbf{本章解决什么问题}：为什么“没事”这两个字，几乎从来不代表没事；它到底返回了什么；以及什么时候该追问，什么时候追问只会把事情弄得更糟。
\end{guideblock}

先给结论：\textbf{“没事”是人类的成功返回码，但它的包体是空的。}空包体本身是一个信号——它不是“没有信号”，而是“有信号，但载荷没发出来”。

值得注意的是，把“没事”当成“没有事”来处理，是人际沟通里代价最高的误读之一。它不会当场报错，它会在三小时后以一条长消息的形式返回。

本章只认领这一句话。全章围绕它做解剖、做改造、做变体。读完的时候，你手里应该有一句能真正用出去的话。

\hcirule

\section{一次 200 空包的返回}

先看样本。

\sample{0610}\par
\begin{mdquote}
你：「你是不是不太高兴？」\par
她：「没事。」\par
你：「那就好。」\par
\vspace{0.3\baselineskip}
三小时之后，你收到一条很长的消息。开头一行是：\par
\vspace{0.3\baselineskip}
「其实我早就想说了。」
\end{mdquote}

拆一下这次交互。

你发出的是一个\textbf{查询请求}：“你是不是不太高兴”。对方返回的是：\textbf{状态码 200，包体为空}。

从协议的角度看，这次请求是“成功”的——对方给了回应，没有报错，对话可以继续。\textbf{但从内容的角度看，你什么都没拿到。}

问题在于，人类接口没有让“成功、本来就没有内容”和“成功、但内容没发出来”区分开的办法。\textbf{对方返回空包体的时候，你收到的只是一个“OK”。}至于是哪一种 OK，要靠别的字段去判断——而大多数人根本不知道还有别的字段。

把这次交互，和一次真的没事放在一起比较。

\sample{0611}\par
\begin{mdquote}
周一早上你出门忘了带伞，淋了一路，到公司楼下头发是湿的。\par
同事看你进来，问：「你怎么了？」\par
你说：「没事，忘带伞了。」\par
他点点头，把纸巾盒推给你，然后你们开始聊今天上午的会。\par
\end{mdquote}

两次交互的状态码是一样的，都是 200，都是“没事”。\textbf{区别不在状态码，在包体。}第二次你附带了内容——“忘带伞了”这四个字就是包体；第一次，那个包的长度是零。

这也解释了为什么“没事”最难处理：\textbf{它是一句正确的话。}对方没有说谎，他确实处在“没事”这个状态里——只是他对“没事”这个词的定义，跟你不一样。你的定义是“没有异常”，他的定义是“不用你管”。

需要说明的是，这不是措辞问题，是接口设计问题。人类这套接口从设计之初就没有给“有事”留一个专用的返回码。所有不便立刻说出口的状态，都被压缩进了同一句“没事”里。\textbf{于是“没事”这个词承担的语义，远远超过了它两个字的容量。}

\section{空包体也带着响应头}

值得注意的是，一个 200 空包体，并不等于一个信息量为零的响应。它仍然带着头（header）。\textbf{人类接口的“没事”，至少带四个头。}真正要读的，是这四个头，而不是包体本身。

\begin{manstable}{P{2.0cm} P{3.8cm} L}
\tbh{响应头} & \tbh{人类接口里的对应} & \tbh{怎么读} \\
\midrule
延迟 & 从你提问到对方回答间隔了多久 & 比平时慢 → 在组织措辞；和你几乎同时出口 → 通常是真的 \\
次数 & 同一句“没事”在这一次对话里出现几次 & 两次以上 → 大概率是希望你再问一次（第三种） \\
信道 & 文字 / 语音 / 电话 / 当面 & 信道越窄，空包体越可疑；当面最快，文字最慢 \\
附随动作 & 说完之后做了什么 & 转身走开 → 第四种；坐回原地 → 第二或第三种；递东西 → 第一种 \\
\bottomrule
\end{manstable}

只看包体，你拿到的是零。看响应头，你至少能拿到一个方向。

\sample{0612}\par
\begin{mdquote}
晚上十一点，你在客厅叠衣服，她从卧室出来接水。你随口问了一句：「今天开会是不是不太顺？」\par
\vspace{0.3\baselineskip}
她接水的动作停了一下，然后说：「没事。」\par
\vspace{0.3\baselineskip}
说完她没有回卧室，坐到你旁边，把剩下的一半衣服接过去叠。\par
\vspace{0.3\baselineskip}
两分钟后她自己开口：「就是那个方案又被打回来了，第三回了。」\par
\end{mdquote}

看这一次的响应头。延迟：接水的动作停了一下，大约一秒——她在决定要不要说。附随动作：她坐下来了，没有走开。\textbf{这两个头指向的是第三种，不是第四种。}你要做的只是别在这个空档里把话题岔开，也别在她开口之前用一串追问把空间填满。

信道这个头，单独值得说。同一句“没事”，信道不同，可信度不同。

\sample{0613}\par
\begin{mdquote}
同一天晚上，你收到两条一模一样的两个字。\par
\vspace{0.3\baselineskip}
朋友 A 发来「没事」，后面跟了一张猫的图。\par
\vspace{0.3\baselineskip}
朋友 B 发来「没事」，三个字，之后再没有输入状态。\par
\vspace{0.3\baselineskip}
A 那条你不用管。B 那条你也不用现在管——但最好在明天中午之前，随便找一件小事跟她说句话。\par
\end{mdquote}

A 和 B 的信道是同一个（文字），差别在附随动作：A 带了图，等于随包体发了一个校验位；B 没带。\textbf{对 B 的正确动作不是追问，而是在第二天发一条跟这件事无关的消息。}它的作用不是解决问题，是证明这条链路还开着。

需要说明的是，四个头的可靠性不同。\textbf{信道和附随动作最可靠，延迟次之，次数最不可靠。}原因很直接：前两个是行为，行为比语言难伪装。

\section{“没事”的四种含义}

值得注意的是，“没事”是一句话，但它至少对应四种完全不同的状态。这四种在字面上无法区分，只能靠上一节那四个响应头，加上第 5 章的方法——先建基线，再看偏离。

\begin{manstable}{P{4.2cm} L L}
\tbh{状态} & \tbh{特征} & \tbh{后续} \\
\midrule
\textbf{一、真的没事} & 语气正常，很快就转到别的话题 & 无 \\
\textbf{二、有事，不想说} & 语气平，回得快，然后这个话题结束了 & 会自己消化，或找别人说 \\
\textbf{三、有事，希望你再问一次} & 语气有保留，回复略慢，可能重复一次“没事” & 你若不问，这条链路关闭 \\
\textbf{四、有事，但希望你别问了} & 语气偏硬，往往伴随其他动作（转身、走开、拿手机） & 你越问，退得越远 \\
\bottomrule
\end{manstable}

如果把这四种含义翻译成机器能懂的状态码，大致是这样——

\begin{manstable}{P{2.6cm} P{2.6cm} L}
\tbh{含义} & \tbh{近似状态码} & \tbh{协议依据} \\
\midrule
一、真的没事 & 204 No Content & 成功，且按设计就没有内容要返回 \\
二、有事，不想说 & 202 Accepted & 已受理，异步处理中，结果稍后另行返回 \\
三、有事，希望你再问一次 & 100 Continue & 临时响应：服务器在等你继续发 \\
四、有事，希望你别问了 & 403 Forbidden & 服务器理解请求但拒绝执行，不应重复请求 \\
\bottomrule
\end{manstable}

需要说明的是，这是一组近似映射，不是严格对应。人类接口没有把话说得这么清楚——\textbf{它把四个状态码压成了同一个 200，然后指望接收方自己去猜。}

\textbf{关键在于第二、三种的差别。}

第二种和第三种在字面上完全一样，都是“没事”。区别只在两个地方：\textbf{回复的速度}，和\textbf{有没有重复}。

\begin{itemize}
\item 回得快、一次说完、然后话题继续 → 大概率是第二种。
\item 回得比平时慢一点、或者说了两次“没事” → 大概率是第三种。
\end{itemize}

需要说明的是，这个判断的准确率并不高。\textbf{在本手册收录的 1,000 段亲密关系对话中，仅凭措辞判断“没事”的含义，准确率约为 58\%。}如果同时读四个响应头，准确率可以提升到约 74\%——这仍然不是一个可以拿来判案的数，但它足以让你知道\textbf{这里可能有东西}。

\sample{0614}\par
\begin{mdquote}
周六下午，你问她要不要去看那场电影。她说：「不用了吧，你忙你的。」\par
\vspace{0.3\baselineskip}
你觉得她语气有点平，多问了一句：「你是不是有点不高兴？」\par
\vspace{0.3\baselineskip}
她说：「没事。」\par
\vspace{0.3\baselineskip}
隔了两秒，又补了一句：「真的没事。」\par
\end{mdquote}

“没事”出现两次，第二次还加了强调。\textbf{强调越用力，越说明第一次不够用。}这是第三种。此时如果收手，对话就结束了；正确的动作是给一条通道、不逼问——比如放下手机说一句“那晚上我早点回来”。\textbf{第三种要的从来不是更多的问题，而是更多的在场。}

\section{追问的问题}

\hciTag{注意事项}：这一节是本手册少数给出明确倾向的地方。

\textbf{默认不要追问。}

原因不是“尊重对方隐私”这种价值判断，而是一个非常实际的计算：\textbf{追问会把“不想说”变成“必须说”。}你成功让对方开口了，但这次的开口是有成本的——对方需要一边整理情绪，一边承受“被逼问”的不适。这个成本记在你的账上，并且它会影响下一次——\textbf{下一次对方会连“没事”都不说，直接不回应。}

追问的代价，是让对方在以后更不愿意对你返回任何状态码。

这不是一个感觉，是可以统计的。\textbf{在 2,000 份涉及追问的对话样本中，一次追问就让对方开口的约占 34\%；而在这部分里，约 71\% 的开口带有明显的“被逼出来”的痕迹}——语速变快、句子变短、不再看你的眼睛。换言之，追问十次，多半只有三四次能问出东西，而问出来的那几次，也未必是对方真的想说的。

\sample{0615}\par
\begin{mdquote}
周二晚上她说“没事”，你追问了四十分钟，最后她把事情说了，也哭了。你以为这次处理得还不错。\par
\vspace{0.3\baselineskip}
周五，她加班到很晚，脸色不好。你问“怎么了”。她说：「没事。」\par
\vspace{0.3\baselineskip}
你说「真的没事？」她看了你一眼，改口说：「有点累，我先睡了。」\par
\end{mdquote}

这不是她变冷淡了。是周二那一次追问，被记在了“追问要付出代价”这一栏里。\textbf{周五她提前给了你一个更省事的答案——“累”。}追问换来的不是这一次的坦白，是下一次的更早封闭。而且你很难察觉，因为表面上看，周二那次你确实“沟通成功”了。

\textbf{但有两种情况必须追问：}

\begin{enumerate}[label=\bfseries\arabic*., leftmargin=2.4em, labelsep=0.7em,
                  itemsep=0.45em, topsep=0.2em, parsep=0.2em]
\item \textbf{涉及安全。}自伤、出走、身体异常、失联。这种情况下，追问与不追问的代价不对等，必须问。
\item \textbf{重复三次以上，且伴随明显的其他偏离。}比如“没事”说了三遍，同时当天没有再回复任何消息。这种情况下，事态已经超过了“不想说”的范围。
\end{enumerate}

需要说明的是，这两种情况都属于例外。例外的意思是：\textbf{平时不这么做，只在满足条件时才做。}如果每次都当成例外来处理，例外本身就失效了。

\section{追问该怎么说}

追问的格式，和大多数人习惯的那句不一样。

\begin{mdquote}
\hciTag{反例}“你到底怎么了？”\par
\vspace{0.3\baselineskip}
\hciTag{反例}“你说没事就肯定有事，说。”\par
\vspace{0.3\baselineskip}
\hciTag{反例}“又来了，你每次都这样。”
\end{mdquote}

这三句的问题在于，它们都是\textbf{要求回答}的请求。

\begin{mdquote}
\hciTag{正例}“行。那我们不说这个了。”\par
\vspace{0.3\baselineskip}
\hciTag{正例}“我在这儿。你想说的时候再说。”\par
\vspace{0.3\baselineskip}
\hciTag{正例}“那我不问了。晚上想吃点什么？”
\end{mdquote}

这三种的结构是：\textbf{确认收到对方的拒绝，然后给出一条不要求回答的通道。}

需要强调的是第三句。\textbf{它做的事情是“把话题换到一个具体的、可以共同处理的问题上”。}人类在情绪上的时候，往往接不住抽象的东西，但接得住“晚上吃什么”。等这顿饭吃完，那句话通常会自己出来。

\sample{0616}\par
\begin{mdquote}
周日中午，他坐在沙发上刷手机，回你话都只回一个字。你问「你是不是不太开心」，他说「没事」。\par
\vspace{0.3\baselineskip}
你没再问。你说：「那我不问了。中午想吃面还是吃饭？」\par
\vspace{0.3\baselineskip}
他愣了一下，说「面吧」。你们下楼吃了面。回来的路上，他自己开口了：「其实是我们组那个项目，今天要交的东西被人改坏了。」\par
\end{mdquote}

关键在第二句“那我不问了”。它做了两件事：第一，确认收到对方的拒绝——我听见了，我不硬来；第二，把话题换到一个具体、可共同处理的问题上。\textbf{追问的失败，多半不是因为问得不够，而是因为问句里只有“你”，没有“我们”。}“你想不想说”没有落点，“中午吃什么”有落点。

\textbf{追问的最高形式，是不追问。}

\section{版本变更记录}

值得注意的是，“没事”这个词的行为，在人类的生命周期里发生过一次静默变更。

\begin{manstable}{P{1.8cm} P{2.4cm} L}
\tbh{版本} & \tbh{时期} & \tbh{“没事”的行为} \\
\midrule
\textbf{v1.0} & 0—3 岁 & “没事”约等于真的没事。此版本尚未具备隐藏状态的能力 \\
\textbf{v2.0} & 12—16 岁 & \textbf{静默变更}。开始出现“没事”作为掩饰。此版本新增“摔门式没事”（说完转身就走） \\
\textbf{v3.0} & 22—35 岁 & “没事”成为常态用语。\textbf{新增“微笑版没事”}——说没事时脸上带着笑，此版本最难解析 \\
\bottomrule
\end{manstable}

\textbf{请注意：此变更不另行通知。}任何一个你熟悉的人，都可能在某个时刻从 v2.0 升级到 v3.0，而你不会收到任何提示——直到某一次你按旧版本的语义去解析，然后发现完全不对。

需要说明的是，上表的年龄区间与附录 B《人类版本变更记录》保持一致。\textbf{“没事”这一项的变更，只是整台机器变更历史里的一小段：v1.0 到 v2.0 之间的那次变更，不只影响“没事”，它是整个接口从直连转为加密的分界点。}

此外还有一个附加说明：\textbf{同一个人的“没事”，针对不同对象，可能运行在不同版本上。}对父母是 v1.0，对伴侣是 v3.0，这是常见的情形。这里的“版本”指的是\textbf{对象轴}，与上表的\textbf{年龄轴}是两个不同的维度——同一个人可以同时处在 v2.0 的年龄，而对伴侣运行 v3.0 的加密策略。

\sample{0617}\par
\begin{mdquote}
周日下午，她妈妈打电话来，问她最近是不是瘦了。她说：「没事妈，就是最近忙。」语气很松，还笑了一下。\par
\vspace{0.3\baselineskip}
挂了电话，你问她：「你妈是不是担心你了？」\par
\vspace{0.3\baselineskip}
她说：「没事。」\par
\vspace{0.3\baselineskip}
同一句“没事”，半小时之内说了两次。第一次是 v1.0，第二次不是。\par
\end{mdquote}

这不是她对你和对她妈两副面孔。这是同一套接口，对不同对象挂了不同的版本。\textbf{对母亲，她的“没事”是一句善意的屏蔽，目的是让对面放心；对你，她的“没事”是一句还没想好怎么解释的占位符。}识别出这一点，比判断她“是不是在骗人”有用得多。

\section{你自己的“没事”该怎么说}

这一节写给说“没事”的那一方。

如果你确实有事，只是现在不想说，那么\textbf{请换一个说法}：

\begin{mdquote}
\hciTag{反例}“没事。”\par
\vspace{0.3\baselineskip}
\hciTag{正例}“我现在不太想说这个。”
\end{mdquote}

这两句话的差别很大。

“没事”是一个\textbf{假的返回码}。它声明“状态正常”，但实际状态不正常。对方收到这个假信号之后，只能自己去猜——而他猜的方向，通常由他自己的不安决定，不由你决定。

“我现在不太想说这个”是一个\textbf{真的返回码}。它明确地告诉对方三件事：第一，确实有事；第二，不是你的问题；第三，这条通道以后还开着。

\textbf{换句话说，“没事”把解析成本转嫁给了对方，而且转嫁得不明不白。}“我现在不想说”把成本落在了自己身上，但它换来的是对方不用瞎猜。

\sample{0618}\par
\begin{mdquote}
你在公司被客户骂了一顿，回到家不想说话。\par
\vspace{0.3\baselineskip}
他问你今天怎么样。你说：「今天不太想说这个，我先缓一会儿。」\par
\vspace{0.3\baselineskip}
他说「行」，去厨房给你倒了杯水，然后打开电视，声音调得很小。\par
\vspace{0.3\baselineskip}
四十分钟后，你自己开口了。\par
\end{mdquote}

这句话只比“没事”多了几个字，但它给了对方三件东西：确认有事、确认不是他的问题、确认通道还开着。\textbf{它的价值不在于解决了问题，而在于把对方从“需要自己猜”的位置上解放了出来。}一个不需要猜的人，才有余地留在这个房间里。

如果你连“有事”都不想承认，还有一句更短的：\textbf{“我缓一缓。”}它同样给出了时间预期，但几乎不暴露内容。它比“没事”诚实，又比“我现在不想说这个”轻。

\section{三个逐字反例}

下面三条都来自真实对话。请逐字看原话，不要先在脑子里把它改成“客气版”再判断——很多话的问题，只在原样的措辞里才看得见。

\textbf{反例一}\par
原话：「我没事，你别管我。」\par
结果：对方照做了。当晚两个人各自睡下，第二天这件事没人再提；一周之后，它成了另一场争吵的开头。\par
替代方案：「我现在心里有点乱，你先让我静一会儿，好了我叫你。」——同样要求空间，但附上了时限和一条通道。

\textbf{反例二}\par
原话：「你说没事就肯定有事，到底说不说。」\par
结果：对方说了，但说的是另一件更小的事。真正的那件事，从此不再对你打开。\par
替代方案：「行，那不说这个。饿不饿？」——收到了拒绝，换一个具体的、可以共同处理的问题。

\textbf{反例三}\par
原话：「你上次也是这样，说没事，结果呢。」\par
结果：对话从“这一次的不高兴”升级成了“你一贯如此”。这一次的“没事”，以后再也不会被打开了。\par
替代方案：不翻旧账，只描述当下：「你今天跟平时不太一样。」——这是观察，不是指控，对方不需要先辩护再开口。

三条反例的共同点，是它们都在\textbf{索要一个对方还没准备好给的返回值}。区别只在于，前两条索要的是内容，第三条索要的是道歉。

\section[常见误区]{常见误区}

\hcimistake{把“没事”直接当成“没有事”}{\hciTag{反例}你问：“你是不是不太高兴？”对方说：“没事。”你答：“哦，那就好。”}{\hciSee{6.1}。这一条是所有误区的源头。三小时后那条长消息，就是这次“哦，那就好”的账单。}

\hcimistake{把追问当成关心的证明}{\hciTag{反例}“你到底怎么了？你不说我怎么帮你？”}{\hciSee{6.4}、6.5。它把“不想说”变成了“必须说”。对方要的不是解决方案，是一条不要求回答的通道。}

\hcimistake{把“没事”理解成拒绝}{\hciTag{反例}对方说“没事”，你回：“那行，我不打扰你了。”}{“没事”通常是“我现在处理不了这件事”，而不是“我不想跟你说话”。这两件事的处理方式完全相反：前者要留出空间，后者要退开距离。用错了，等于对方刚关上一扇门，你把整条走廊也封了。}

\hcimistake{自己说“没事”，然后怪对方没看出来}{\hciTag{反例}你：“没事。”……（三小时后）“你根本不在乎我，我说没事你就真信了。”}{\hciSee{6.7}。你发的是空包，却期待对方解析出内容。对方连包体都没收到，不可能解析出连你自己都没写下来的东西。}

\hcimistake{隔三天再翻旧账}{\hciTag{反例}“那天你就说没事，其实……”}{“其实”两个字一出来，对方的当下和你的当下被同时否定——他当时不想说，你现在才发作，两件事被并成一笔账。}

\hcimistake{把对方的“没事”转述给别人}{\hciTag{反例}“她跟我说没事，但我觉得她肯定有事。”}{这句话一旦传回她耳朵里，损害是永久性的。你不是在求解，你是在替她给第三方发布一个她拒绝发布的返回码。}

\hcimistake{用一串追问代替等待}{\hciTag{反例}“怎么了？”“真的没事？”“你确定？”“那你到底为什么不高兴？”}{四句连发，等于把服务器的“请稍后重试”当成了“立即重试”。重试次数越多，服务端只会把超时时间调得越长。\hciSee{6.2}。}

\hcimistake{在对方说“没事”之后，立刻开始证明自己没事}{\hciTag{反例}对方说“没事”，你开始讲道理：“我明明没有那个意思，那天我……”}{对方要的不是道理，是那条通道。你越急着自证清白，越是在把话题拉回到“谁对谁错”，而对方此刻根本没打算审理这件事。}

\section{延伸：当对方是长辈、上级或陌生人}

同一句“没事”，换一个关系，读法和做法都要换。这一节处理三种关系。

\textbf{当对方是长辈}：长辈的“没事”通常是第二种或第四种，而且他们更少主动说出“希望你再问一次”这一种。对他们，追问往往无效，甚至反作用。\textbf{不要追问，改用行动。}倒杯水、坐下、不提这件事，比任何话术都有用。

\sample{0619}\par
\begin{mdquote}
你妈住院复查回来，坐在沙发上不说话。你问她结果怎么样。她说：「没事，老毛病。」\par
\vspace{0.3\baselineskip}
你没再问，去厨房倒了杯温水放在她手边，然后坐在旁边看手机。\par
\vspace{0.3\baselineskip}
过了一会儿，她自己说：「医生说那个指标比上回高了一点，让下周再去看看。」\par
\end{mdquote}

这里的动作不是客套。\textbf{倒水、坐下、留在同一个房间里，本身就是一条不要求回答的通道。}长辈不擅长收下一条抽象的“我在”，但收得下一杯递到手边的水。

\textbf{当对方是上级}：上级说“没事”，多数时候确实是公事上的没事。\textbf{此时追问是减分的}——它会被理解为你在探查，而不是在关心。正确做法是记下来，等下一次汇报时用结果去回应。

\begin{mdquote}
你：「张总，那个客户的反馈，您看要不要先把方案调一下？」\par
他：「没事，先这样。」\par
你说：「好，那我记一下，下次汇报的时候把改动一并给您看。」\par
\end{mdquote}

最后一句是关键。它没有追问，但做了一件更划算的事：\textbf{把它变成了一个后续要交付的项。}上级的“没事”如果是真的，你多记一笔没有成本；如果是假的，你在下一次就把台阶递到了。

\textbf{当对方是陌生人}：陌生人的“没事”通常是真的没事。\textbf{不加解读。}把处理亲密关系的规则硬套到陌生人身上——反复追问“你真的没事吗”——会被读成越界，而不是关心。关系的带宽越窄，能承载的解读就越少。

\hcirule

%% ---------- 本章小结 ----------
\hciblocktitle{bullseye}{本章小结}

综上所述，本章的核心结论可以概括为以下几点：

\begin{enumerate}[label=\bfseries\arabic*., leftmargin=2.4em, labelsep=0.7em,
                  itemsep=0.45em, topsep=0.2em, parsep=0.2em]
\item “没事”是 200 状态码，但包体为空。空包体是信号，不是没有信号。
\item 空包体也带着响应头：延迟、次数、信道、附随动作。四个头合起来读，比包体本身更有信息。
\item “没事”至少有四种含义，字面上无法区分；这四种可以近似映射到 204、202、100、403 四个状态码。
\item 默认不追问。追问会把“不想说”变成“必须说”，代价记在下一次。
\item 追问的最高形式是不追问——换一个具体的、可以共同处理的问题。
\item 你自己不想说的时候，别说“没事”，说“我现在不想说这个”，或者退一步说“我缓一缓”。
\item 关系不同，读法不同：对长辈改用行动，对上级改用记录，对陌生人不加解读。
\end{enumerate}

%% ---------- 练习 ----------
\begin{exercisessec}
\item 回忆你最近一次听到“没事”的场景。判断它属于四种含义里的哪一种，并写下你当时是怎么处理的。（[请在此处填写当时的情境]）

\item 在下一次对方说“没事”的时候，不追问，改用一个具体问题（比如“晚上吃什么”）。记录后来那句话有没有自己出来。

\item 把你自己最近说过的一次“没事”调出来，改写成一个真实的返回码。写下改后的原话。

\item 找一个你最近一次说“没事”的场合，把当时能读到的四个响应头列出来——延迟、次数、信道、附随动作——然后判断：对方当时读到的，最可能是哪一种含义？和你实际的意思差了多少？

\item 找一位长辈，在一件他没有主动说的事上，刻意用一次“动作代替追问”。记录他后来有没有自己开口，用了多久。
\end{exercisessec}

%% ---------- 自测题 ----------
\begin{quizsec}
\item 对方说“没事”时，下列哪一种做法是正确的？

\begin{enumerate}[label=\Alph*., leftmargin=2.2em, labelsep=0.6em,
                  itemsep=0.2em, topsep=0.3em, parsep=0.1em]
\item 立刻当成没事，转移话题
\item 追问“你到底怎么了”，直到对方说
\item 确认收到，给一条不要求回答的通道
\item 把这件事记下来，下次翻出来
\end{enumerate}

\item 判断：追问得越紧，对方越会觉得你关心他。

\item 简答：为什么“我现在不想说这个”比“没事”更好？

\item 以下哪一项属于“附随动作”这个响应头？

\begin{enumerate}[label=\Alph*., leftmargin=2.2em, labelsep=0.6em,
                  itemsep=0.2em, topsep=0.3em, parsep=0.1em]
\item 对方回复所用的时间
\item 对方说完之后转身走开
\item “没事”在对话里出现的次数
\item 消息是文字还是语音
\end{enumerate}

\item 简答：对长辈说“没事”时，为什么本手册建议用行动代替追问？
\end{quizsec}

%% ---------- 参考答案 ----------
\begin{answerssec}
\item C。“没事”是空包体，既不等于没事，也不等于拒绝。正确做法是确认收到对方的拒绝，然后给出一条不要求回答的通道。\hciSee{6.5}。

\item 错。追问会把“不想说”变成“必须说”，这次的开口是有成本的，而这个成本会影响下一次——下一次对方会连“没事”都不说。\hciSee{6.4}。

\item 因为“没事”是一个假的返回码，它声明状态正常但实际不正常，把解析成本转嫁给了对方，而且转嫁得不明不白。“我现在不想说这个”是真话，它同时说明：确实有事、不是你的问题、通道以后还开着。\hciSee{6.7}。

\item B。附随动作指的是说完之后做了什么——转身走开、坐回原地、递东西。回复所用的时间属于“延迟”，出现次数属于“次数”，文字或语音属于“信道”。\hciSee{6.2}。

\item 因为长辈的“没事”多为第二种或第四种，他们比年轻人更少主动希望被追问。对长辈，追问往往无效甚至反作用；而倒水、坐下这类动作本身就是一条不要求回答的通道，他们收得下动作，收不下抽象的话。\hciSee{6.10}。
\end{answerssec}

\hcirule

\begin{qabox}
\qaQ{读者提问}{如果对方永远都在说“没事”，从来不直说，我是不是只能一直猜？}
\qaA{本手册回答}{这是一个非常好的问题。需要说明的是，“永远说没事”通常不是沟通技巧问题，而是这段关系里“直说”是否安全的问题。本手册能教的是怎么正确地接收，不能替一段关系解决安全感。需要区分的是：偶尔的“没事”是接口的常态，长期的“没事”是一个诊断信号——它说明对方在你这里，没有找到过一条“直说之后不会更糟”的通道。}
\end{qabox}

\hcirule

希望本章的内容能对你有所帮助。如需进一步了解第 7 章《反话》的相关内容，我可以随时为你展开。

%% 章末「页脚交互条」由 hci-manual.cls 自动挂接，无需手写。
